\documentclass[../../main.tex]{subfiles}
\begin{document}
\chapter{Algorithms and computation}
\label{ch:computation}

\section{Pseudocode}
The Euclidean algorithm repeatedly replaces a pair by a smaller pair with
the same greatest common divisor. It is a useful example of mathematical
reasoning alongside pseudocode. For broader background on discrete methods,
see \textcite{concrete1994}.

\begin{algorithm}[htbp]
  \caption{Euclidean algorithm}
  \label{alg:euclid}
  \begin{algorithmic}[1]
    \Require Integers $a\ge b\ge0$ with $a>0$
    \Ensure $\gcd(a,b)$
    \While{$b\ne0$}
      \State $r\gets a\bmod b$
      \State $a\gets b$
      \State $b\gets r$
    \EndWhile
    \State \Return $a$
  \end{algorithmic}
\end{algorithm}

\begin{prop}\label{prop:gcd-invariant}
Let $a\ge b>0$ be integers. If $a=qb+r$ with $q\in\ZZ$ and
$0\le r<b$, then $\gcd(a,b)=\gcd(b,r)$.
\end{prop}
\begin{proof}
A positive integer divides both $a$ and $b$ if and only if it divides both
$b$ and $r=a-qb$. Thus the two pairs have the same common positive divisors
and the same greatest common divisor.
\end{proof}

Each iteration of \cref{alg:euclid} preserves the greatest common divisor
by \cref{prop:gcd-invariant}. The new second entry satisfies $0\le r<b$, so
the algorithm terminates. At termination, $\gcd(a,0)=a$.

\clearpage
\section{A matching implementation}
The \texttt{listings} package typesets source code directly, without an
external syntax-highlighting program or shell escape. This template defines
the \texttt{everett-code} style in the preamble.

\begin{lstlisting}[language=Python,caption={Computing a greatest common divisor},label={lst:gcd}]
def gcd(a, b):
    """Return gcd(a, b) for nonnegative integers, not both zero."""
    while b:
        a, b = b, a % b
    return a

assert gcd(252, 105) == 21
assert gcd(17, 0) == 17
\end{lstlisting}

The simultaneous assignment keeps the old values of \texttt{a} and
\texttt{b} while computing the new pair. For $(252,105)$, the successive
pairs are shown across the main column and margin:
\begin{fullwidth}
\[
  (252,105)\taking{}(105,42)\taking{}(42,21)\taking{}(21,0).
\]
\end{fullwidth}
When showing a longer program, store it in its own file and include it with:
\begin{lstlisting}[language={[LaTeX]TeX}]
\lstinputlisting[language=Python]{path/to/file.py}
\end{lstlisting}

\section{Inline lists and short solutions}
The \texttt{parlist} environment is useful for a compact checklist. For
example, a proof of an algorithm should establish
\begin{parlist}
  \item the preconditions,
  \item an invariant,
  \item termination, and
  \item the postcondition.
\end{parlist}

\begin{exercise}[An integer pair]\label{ex:gcd-lcm}
Compute $\gcd(84,30)$ and $\lcm(84,30)$.
\end{exercise}
\begin{soln}
The Euclidean algorithm gives $\gcd(84,30)=6$. Since the inputs are
positive, their least common multiple is
$\lcm(84,30)=84\cdot30/6=420$.
\end{soln}

\begin{recap}
An algorithm needs both an invariant and a reason to terminate.
The Euclidean algorithm preserves the common divisors while strictly
decreasing its nonzero second input. The code mirrors those two ideas.
\end{recap}
\end{document}
